\documentclass[11pt]{article}
\usepackage[margin=.8in,top=.7in,bottom=.7in]{geometry}
\usepackage{amsmath,amssymb,amsthm,hyperref}
\renewcommand{\normalsize}{\fontsize{10}{12}\selectfont}
\DeclareMathSizes{10}{10}{8}{6}
\hypersetup{colorlinks=true,linkcolor=blue,urlcolor=blue,citecolor=blue}
\setlength{\emergencystretch}{2em}
\newtheorem{theorem}{Theorem}
\newcommand{\I}{I_2}
\title{A cubic counterexample to three\\C*-algebraic polynomial inequalities}
\author{Research proof and formalization packet}
\date{October 26, 2026}
\begin{document}
\begin{center}
{\Large\bfseries A cubic counterexample to three\\
C*-algebraic polynomial inequalities}\\[.5em]
Research proof and formalization packet\\
October 26, 2026
\end{center}
\begin{abstract}
We give one degree-three matrix construction refuting Krishna's
C*-algebraic Schoenberg, de Bruin--Sharma and Kushel--Tyaglov
conjectures. Its ordered deletion derivative factors as required for
every argument in $M_2(\mathbb C)$, including arguments that do not
commute with its roots. Both adjoint versions of each proposed
inequality fail, with exact negative quadratic forms. A complete
36-declaration Lean certificate proves the factorization and all six
Hermitian positivity failures. The classical scalar theorems are
unaffected.
\end{abstract}

\section{The source assertions}
Krishna \cite{K}, Section~2, considers
$P(z)=(z-a_1)\cdots(z-a_d)$ over a C*-algebra $\mathcal A$ and defines
\[
 P'(z)=\sum_{j=1}^d(z-a_1)\cdots\widehat{(z-a_j)}\cdots(z-a_d),
 \qquad z\in\mathcal A.
\]
The factors retain their order. The hypotheses of his Conjectures~2.1,
2.3 and~2.4 require a factorization
$P'(z)=d(z-b_1)\cdots(z-b_{d-1})$ on the whole algebra.
They impose no normality or self-adjointness condition on the $a_j$.
MathDB~355008 and~355009 record the first two conjectures.

For centered roots $\sum a_j=0$ and degree $d=3$, write
\[
 S=\sum_j a_ja_j^*,\quad T=\sum_j(a_ja_j^*)^2,\quad
 C=\sum_j a_j^2,\quad E=\sum_j a_j^2(a_j^2)^*.
\]
The left-adjoint versions of the three assertions specialize to
\begin{align}
 \sum_{k=1}^2 b_kb_k^*&\preceq S/3,\label{S}\\
 \sum_{k=1}^2(b_kb_k^*)^2&\preceq 2S^2/9-T/3,\label{D}\\
 \sum_{k=1}^2(b_kb_k^*)^2&\preceq -T+S^2/9+CC^*/9+2E/3.\label{KT}
\end{align}
Each has a corresponding adjoint-reversed assertion: replace
$a_ja_j^*$ and $b_kb_k^*$ by $a_j^*a_j$ and $b_k^*b_k$,
$CC^*$ by $C^*C$, and $E$ by $\sum_j(a_j^2)^*a_j^2$.
In the third assertion the terms involving $\sum a_j$ vanish.
We take $n=d=3$, so the source's mixed $n,d$ indexing in
Conjecture~2.4 has no bearing on this example.

The scalar predecessors are known theorems; see
Cheung--Ng~\cite{CN} and Kushel--Tyaglov~\cite{KT}.
The result below concerns their proposed C*-algebraic extensions.

\section{One matrix construction}
\begin{theorem}\label{main}
All three C*-algebraic conjectures fail at degree $3$ in
$\mathcal A=M_2(\mathbb C)$. One centered cubic, with a derivative
factorization valid for every $z\in\mathcal A$, violates both adjoint
versions of all three inequalities.
\end{theorem}
\begin{proof}
Put
\[
 R=\begin{pmatrix}1&2\\0&-1\end{pmatrix},\qquad
 (a_1,a_2,a_3)=(R,-2R,R),\qquad (b_1,b_2)=(\I,-\I).
\]
Then $R^2=\I$ and $\sum a_j=0$. Without interchanging $z$ and $R$,
the ordered derivative expands to
\begin{align*}
 P'(z)&=(z+2R)(z-R)+(z-R)^2+(z-R)(z+2R)\\
      &=3z^2-3R^2=3(z-\I)(z+\I).
\end{align*}
Thus the full factorization hypothesis holds, even at noncommuting
arguments. Each $a_j$ is a zero of the displayed cubic and each $b_k$
is a zero of its deletion derivative.

Let
\[
 H=RR^*=\begin{pmatrix}5&-2\\-2&1\end{pmatrix},\qquad
 J=R^*R=\begin{pmatrix}1&2\\2&5\end{pmatrix}.
\]
We have $S=6H$, $T=18H^2$, $C=6\I$ and $E=18\I$.
Both critical-point sums in \eqref{S}--\eqref{KT} equal $2\I$.
Consequently, the right sides of those inequalities are respectively
\[
 2H,\qquad 2H^2,\qquad 16\I-14H^2.
\]
The reversed versions have the same expressions with $H$ replaced by $J$.

For $v=(1,2)^T$ and $w=(2,-1)^T$,
$Hv=(1,0)^T$, $Jw=(0,-1)^T$ and $v^*v=w^*w=5$.
It follows that
\[
 v^*2(H-\I)v=v^*2(H^2-\I)v=-8,
 \qquad
 w^*2(J-\I)w=w^*2(J^2-\I)w=-8.
\]
These refute the first two pairs of inequalities. Finally,
$ (H^2)_{11}=(J^2)_{22}=29 $, so
\[
 e_1^*\bigl(14\I-14H^2\bigr)e_1
 =e_2^*\bigl(14\I-14J^2\bigr)e_2=-392.
\]
These are the defects in the third pair. All six defects are
real symmetric, hence Hermitian, and none is positive semidefinite.
\end{proof}

The roots even commute pairwise. Their nonnormality is essential to
this construction: similarity that is not unitary does not preserve
the adjoint/order expressions appearing in these conjectures.
No new theorem about nonnormal involutions is being asserted.

\section{Formal scope and provenance}
The two Lean modules construct the literal ordered cubic and deletion
derivative, prove its factorization for every complex $2\times2$ matrix,
and prove centeredness, the actual root/critical-point evaluations,
all six exact defect matrices, Hermitianity and failure of positive
semidefiniteness. The final closed theorem is
\texttt{CstarPolynomial.source\_all\_three\_conjectures\_false}.
In particular, the full terms of Conjecture~2.4 are formalized before
their centered specialization; they are not assumed to vanish.

Strict replay uses Lean~4.34.1 and mathlib revision
\texttt{d13f23b723b8a846827a245b89c10fc7d3f11612},
with exact dependency pins included. Every declaration type and its
axioms are recorded; only \texttt{propext}, \texttt{Classical.choice}
and \texttt{Quot.sound} occur. A separate rational-arithmetic script
checks the six expressions, while Lean supplies the universal proof.
Acceptance and the status of independent review and bounded literature
searches are recorded separately. No absolute priority claim is made.
The construction and its related failures constitute one outcome.

\begin{thebibliography}{9}
\bibitem{K} K.~Mahesh Krishna,
\emph{C*-algebraic Schoenberg Conjecture},
arXiv:2206.06653v1 (2022), Conjectures~2.1, 2.3 and~2.4.
\bibitem{CN} W.~S.~Cheung and T.~W.~Ng,
\emph{A companion matrix approach to the study of zeros and critical
points of a polynomial}, J.~Math.~Anal.~Appl. \textbf{319} (2006), 690--707.
\bibitem{KT} O.~Kushel and M.~Tyaglov,
\emph{Circulants and critical points of polynomials},
J.~Math.~Anal.~Appl. \textbf{439} (2016), 634--650.
\end{thebibliography}
\end{document}
